\documentclass[10pt,a4paper]{amsart}
\usepackage[margin=19mm]{geometry}
\usepackage{amsmath,amssymb,mathtools}
\usepackage[scr=rsfs,cal=euler]{mathalfa}
\usepackage{xfrac}
\usepackage[unicode,hidelinks,bookmarks=false]{hyperref}
\newcommand{\ie}{i.e.\ }
\newcommand{\cf}{cf.\ }
\newcommand{\resp}{respectively\ }
\newcommand{\Subsection}{\subsection}
\providecommand{\nobreakdash}{\nobreak\hbox{-}}
\setlength{\emergencystretch}{2em}
\multlinegap0pt
\usepackage{bm}
\usepackage{mdframed}
\usepackage{enumerate}
\usepackage[shortlabels]{enumitem}
\usepackage{bbm}
\usepackage{dsfont}
\usepackage{amssymb}
\usepackage{mathtools}
\usepackage{xcolor}
\usepackage{cancel}
\usepackage{xfrac}
\usepackage{siunitx}
\usepackage{float}
\usepackage[all]{xy}
\usepackage{tikz}
\usepackage{booktabs}
\usepackage{multirow}
\usepackage{natbib}
\usepackage{stackengine}
\newtheorem{theorem}{Theorem}
\def\P{{\mathrm P}}
\def\d{{\mathrm d}}
\def\given{\,|\,}
\newcommand{\htau}{\widehat\tau}
\newcommand{\indep}{\perp \!\!\!\perp}
\title{\LARGE Introducing the b-value: combining unbiased and biased estimators from a sensitivity analysis perspective}
\author{Zhexiao Lin, Peter J. Bickel, Peng Ding}
\date{Source: arXiv:2602.16310v1 (CC BY 4.0)}
\begin{document}
\maketitle

\noindent\emph{Source and licence.} \LARGE Introducing the b-value: combining unbiased and biased estimators from a sensitivity analysis perspective. Version arXiv:2602.16310v1 (CC BY 4.0), distributed by arXiv under the Creative Commons Attribution 4.0 International licence, http://creativecommons.org/licenses/by/4.0/, as recorded in the arXiv licence metadata for this exact version and shown on the abstract page https://arxiv.org/abs/2602.16310v1. Full licence notice: \texttt{licenses/arxiv-2602.16310-CC-BY-4.0.txt}.\par\smallskip

\noindent\textit{Editorial scope.} Counted unit: the article's theorem prop:CI\_multi\_ST (source0.tex lines 820--832). The article's complete original proof of it is delivered with it (source0.tex lines 2009--2106, one \texttt{proof} environment of 98 lines, which the article places away from the statement). No mathematics is written, repaired or re-derived: the statement and the proof are the article's own lines, in the article's own notation, and no retained line is substituted or re-flowed.  No further source-local context is needed: every cross-reference of every retained line resolves inside the two delivered spans.  Nothing was omitted from any retained span, so the package carries no omission record.  No retained line contains a citation, so the wrapper carries no bibliography and no bibliography entry.  No retained line contains a figure environment, an image-inclusion command or drawing code, so no image is delivered and the package declares no asset dependency.  Editorial formatting only, in addition: an AMS \texttt{amsart} preamble that restates the article's own declarations and macros used by the retained lines, namely the environments \texttt{theorem} on separate counters, together with the standard packages those lines need.  The retained environments print in this document as Theorem 1 in order of appearance, whereas the article prints the same items with the numbering of its own full document.  The complete native source of the article as distributed in the arXiv e-print for this exact version is bundled unchanged as \texttt{sources/194869/source0.tex}.
\begin{theorem}\label{prop:CI_multi_ST}
  $\hat{M}_{\rm ST}$ is the solution to the equation of $M$:
  \begin{align*}
    \inf_{\bm{\Delta}: \bm{\Sigma}^{-1/2} \bm{\Delta} = \bm{b} \odot \bm{s}, \bm{s} \in \{-1,1\}^d} \P_{\bm{\Delta}}((\bm{\htau}_{\rm ST} - \bm{\tau})^\top (\bm{\Sigma}_0^{-1} + \bm{\Sigma}_1^{-1}) (\bm{\htau}_{\rm ST} - \bm{\tau}) \le M) = 1-\zeta,
  \end{align*} 
  with an explicit form given by:
  \begin{align}\label{eq:P_multi_ST}
    &\P_{\bm{\Sigma}^{-1/2}\bm{\Delta} = \bm{t}}\left((\bm{\htau}_{\rm ST}-\bm{\tau})^\top (\bm{\Sigma}_0^{-1} + \bm{\Sigma}_1^{-1})(\bm{\htau}_{\rm ST}-\bm{\tau})\le M\right) \nonumber\\
    =& \Psi_d\left(M;\left\lVert (\bm{\Sigma}_0^{-1} + \bm{\Sigma}_1^{-1})^{-1/2} \bm{\Sigma}_1^{-1} \bm{\Sigma}^{1/2} \bm{t}\right\rVert_2^2\right) \Psi_d\left(q; \left\lVert (\bm{\Sigma}_0 + \bm{\Sigma}_1)^{-1/2} \bm{\Sigma}^{1/2} \bm{t}\right\rVert_2^2\right) \nonumber\\
    &+ \int_{\lVert \bm{u} \rVert_2^2 > q} \Psi_d\left(M;\left\lVert (\bm{\Sigma}_0^{-1} + \bm{\Sigma}_1^{-1})^{-1/2} \bm{\Sigma}_1^{-1} [\bm{\Sigma}^{1/2} \bm{t} - (1-h_q^*(\lVert \bm{u} \rVert_2^2))(\bm{\Sigma}_0+ \bm{\Sigma}_1)^{1/2} \bm{u}]\right\rVert_2^2\right) \nonumber\\
    &\phi_{(\bm{\Sigma}_0 + \bm{\Sigma}_1)^{-1/2} \bm{\Sigma}^{1/2} \bm{t}, \bm{I}_d}\left(\bm{u}\right) \d \bm{u}.
  \end{align}
\end{theorem}

\begin{proof}[Proof of Theorem~\ref{prop:CI_multi_ST}]
Recall that the soft-thresholding estimator $\bm{\htau}_{\rm ST}$ is defined as
\[
  \bm{\hat\tau}_{\rm ST} = \bm{\htau}_0 + h_q(\lVert (\bm{\Sigma}_0+ \bm{\Sigma}_1)^{-1/2}(\bm{\htau}_1 - \bm{\htau}_0) \rVert_2^2)(\bm{\Sigma}_0^{-1} + \bm{\Sigma}_1^{-1})^{-1}\bm{\Sigma}_1^{-1}(\bm{\htau}_1 - \bm{\htau}_0).
\]
A direct calculation shows that
\begin{align*}
  \bm{\htau}_0 + (\bm{\Sigma}_0^{-1} + \bm{\Sigma}_1^{-1})^{-1} \bm{\Sigma}_1^{-1} (\bm{\htau}_1 - \bm{\htau}_0) \indep \bm{\htau}_1 - \bm{\htau}_0.
\end{align*}
Moreover,
\[
  \bm{\htau}_1 - \bm{\htau}_0 \sim N(\bm{\Delta}, \bm{\Sigma}_0+ \bm{\Sigma}_1),
\]
and
\[
  \bm{\htau}_0 + (\bm{\Sigma}_0^{-1} + \bm{\Sigma}_1^{-1})^{-1} \bm{\Sigma}_1^{-1} (\bm{\htau}_1 - \bm{\htau}_0) \sim N(\bm{\tau} + (\bm{\Sigma}_0^{-1} + \bm{\Sigma}_1^{-1})^{-1} \bm{\Sigma}_1^{-1} \bm{\Delta}, (\bm{\Sigma}_0^{-1} + \bm{\Sigma}_1^{-1})^{-1} ).
\]

We first consider the event that the pretest accepts, $\lVert (\bm{\Sigma}_0+ \bm{\Sigma}_1)^{-1/2}(\bm{\htau}_1 - \bm{\htau}_0) \rVert_2^2 \le q$. Conditional on this event, $\bm{\htau}_{\rm ST} = \bm{\htau}_0 + (\bm{\Sigma}_0^{-1} + \bm{\Sigma}_1^{-1})^{-1} \bm{\Sigma}_1^{-1} (\bm{\htau}_1 - \bm{\htau}_0) $, and therefore
\begin{align*}
  (\bm{\htau}_{\rm ST} - \bm{\tau}) \given \{\lVert (\bm{\Sigma}_0+ \bm{\Sigma}_1)^{-1/2}(\bm{\htau}_1 - \bm{\htau}_0) \rVert_2^2 \le q\} \sim N((\bm{\Sigma}_0^{-1} + \bm{\Sigma}_1^{-1})^{-1} \bm{\Sigma}_1^{-1} \bm{\Delta}, (\bm{\Sigma}_0^{-1} + \bm{\Sigma}_1^{-1})^{-1} ).
\end{align*}
After standardization,
\begin{align*}
  &(\bm{\Sigma}_0^{-1} + \bm{\Sigma}_1^{-1})^{1/2} (\bm{\htau}_{\rm ST} - \bm{\tau}) \given \{\lVert (\bm{\Sigma}_0+ \bm{\Sigma}_1)^{-1/2}(\bm{\htau}_1 - \bm{\htau}_0) \rVert_2^2 \le q\} \\
  \sim& N((\bm{\Sigma}_0^{-1} + \bm{\Sigma}_1^{-1})^{-1/2} \bm{\Sigma}_1^{-1} \bm{\Delta}, \bm{I}_d),
\end{align*}
which implies
\begin{align*}
  &(\bm{\htau}_{\rm ST} - \bm{\tau})^\top (\bm{\Sigma}_0^{-1} + \bm{\Sigma}_1^{-1}) (\bm{\htau}_{\rm ST} - \bm{\tau}) \given \{\lVert (\bm{\Sigma}_0+ \bm{\Sigma}_1)^{-1/2}(\bm{\htau}_1 - \bm{\htau}_0) \rVert_2^2 \le q\} \\
  \sim& \chi^2_d\left(\left\lVert (\bm{\Sigma}_0^{-1} + \bm{\Sigma}_1^{-1})^{-1/2} \bm{\Sigma}_1^{-1} \bm{\Delta}\right\rVert_2^2\right).
\end{align*}

We next consider the event that the pretest rejects, $\lVert (\bm{\Sigma}_0+ \bm{\Sigma}_1)^{-1/2}(\bm{\htau}_1 - \bm{\htau}_0) \rVert_2^2 > q$. Conditional on $(\bm{\Sigma}_0+ \bm{\Sigma}_1)^{-1/2}(\bm{\htau}_1 - \bm{\htau}_0) = \bm{u}$, we have
\[
  (\bm{\Sigma}_0^{-1} + \bm{\Sigma}_1^{-1})^{-1} \bm{\Sigma}_1^{-1} (\bm{\htau}_1 - \bm{\htau}_0) = (\bm{\Sigma}_0^{-1} + \bm{\Sigma}_1^{-1})^{-1} \bm{\Sigma}_1^{-1} (\bm{\Sigma}_0+ \bm{\Sigma}_1)^{1/2} \bm{u}.
\]
Since $\bm{\htau}_{\rm ST} = \bm{\htau}_0 + h_q^*(\lVert (\bm{\Sigma}_0+ \bm{\Sigma}_1)^{-1/2}(\bm{\htau}_1 - \bm{\htau}_0) \rVert_2^2)(\bm{\Sigma}_0^{-1} + \bm{\Sigma}_1^{-1})^{-1}\bm{\Sigma}_1^{-1}(\bm{\htau}_1 - \bm{\htau}_0)$ conditional on this event, we have
\begin{align*}
  &(\bm{\htau}_{\rm ST} - \bm{\tau}) \given \{(\bm{\Sigma}_0+ \bm{\Sigma}_1)^{-1/2}(\bm{\htau}_1 - \bm{\htau}_0) = \bm{u}, \lVert (\bm{\Sigma}_0+ \bm{\Sigma}_1)^{-1/2}(\bm{\htau}_1 - \bm{\htau}_0) \rVert_2^2 > q\} \\
  =& \bm{\htau}_0 + h_q^*(\lVert \bm{u} \rVert_2^2)(\bm{\Sigma}_0^{-1} + \bm{\Sigma}_1^{-1})^{-1}\bm{\Sigma}_1^{-1}(\bm{\htau}_1 - \bm{\htau}_0) \given \{(\bm{\Sigma}_0+ \bm{\Sigma}_1)^{-1/2}(\bm{\htau}_1 - \bm{\htau}_0) = \bm{u}, \lVert (\bm{\Sigma}_0+ \bm{\Sigma}_1)^{-1/2}(\bm{\htau}_1 - \bm{\htau}_0) \rVert_2^2 > q\}\\
  =& \bm{\htau}_0 + (\bm{\Sigma}_0^{-1} + \bm{\Sigma}_1^{-1})^{-1} \bm{\Sigma}_1^{-1} (\bm{\htau}_1 - \bm{\htau}_0) \given \{(\bm{\Sigma}_0+ \bm{\Sigma}_1)^{-1/2}(\bm{\htau}_1 - \bm{\htau}_0) = \bm{u}, \lVert (\bm{\Sigma}_0+ \bm{\Sigma}_1)^{-1/2}(\bm{\htau}_1 - \bm{\htau}_0) \rVert_2^2 > q\}\\
  &- (1-h_q^*(\lVert \bm{u} \rVert_2^2))(\bm{\Sigma}_0^{-1} + \bm{\Sigma}_1^{-1})^{-1} \bm{\Sigma}_1^{-1} (\bm{\Sigma}_0+ \bm{\Sigma}_1)^{1/2} \bm{u}\\
  \sim& N((\bm{\Sigma}_0^{-1} + \bm{\Sigma}_1^{-1})^{-1} \bm{\Sigma}_1^{-1} [\bm{\Delta} - (1-h_q^*(\lVert \bm{u} \rVert_2^2))(\bm{\Sigma}_0+ \bm{\Sigma}_1)^{1/2} \bm{u}], (\bm{\Sigma}_0^{-1} + \bm{\Sigma}_1^{-1})^{-1} ).
\end{align*}
After standardization,
\begin{align*}
  &(\bm{\Sigma}_0^{-1} + \bm{\Sigma}_1^{-1})^{1/2} (\bm{\htau}_{\rm ST} - \bm{\tau}) \given \{(\bm{\Sigma}_0+ \bm{\Sigma}_1)^{-1/2}(\bm{\htau}_1 - \bm{\htau}_0) = \bm{u}, \lVert (\bm{\Sigma}_0+ \bm{\Sigma}_1)^{-1/2}(\bm{\htau}_1 - \bm{\htau}_0) \rVert_2^2 > q\} \\
  \sim& N((\bm{\Sigma}_0^{-1} + \bm{\Sigma}_1^{-1})^{-1/2} \bm{\Sigma}_1^{-1} [\bm{\Delta} - (1-h_q^*(\lVert \bm{u} \rVert_2^2))(\bm{\Sigma}_0+ \bm{\Sigma}_1)^{1/2} \bm{u}], \bm{I}_d),
\end{align*}
which implies
\begin{align*}
  &(\bm{\htau}_{\rm ST} - \bm{\tau})^\top (\bm{\Sigma}_0^{-1} + \bm{\Sigma}_1^{-1}) (\bm{\htau}_{\rm ST} - \bm{\tau}) \given \{(\bm{\Sigma}_0+ \bm{\Sigma}_1)^{-1/2}(\bm{\htau}_1 - \bm{\htau}_0) = \bm{u}, \lVert (\bm{\Sigma}_0+ \bm{\Sigma}_1)^{-1/2}(\bm{\htau}_1 - \bm{\htau}_0) \rVert_2^2 > q\} \\
  \sim& \chi^2_d\left(\left\lVert (\bm{\Sigma}_0^{-1} + \bm{\Sigma}_1^{-1})^{-1/2} \bm{\Sigma}_1^{-1} [\bm{\Delta} - (1-h_q^*(\lVert \bm{u} \rVert_2^2))(\bm{\Sigma}_0+ \bm{\Sigma}_1)^{1/2} \bm{u}]\right\rVert_2^2\right).
\end{align*}

We now evaluate the worst-case coverage probability of the confidence region
\[
  \{\bm{\tau}\in\mathbb{R}^d:(\bm{\htau}_{\rm ST}-\bm{\tau})^\top (\bm{\Sigma}_0^{-1} + \bm{\Sigma}_1^{-1})(\bm{\htau}_{\rm ST}-\bm{\tau})\le M\}.
\]
We decompose the coverage probability according to whether the
pretest accepts or rejects:
\begin{align*}
  &\P_{\bm{\Delta}}\left((\bm{\htau}_{\rm ST}-\bm{\tau})^\top (\bm{\Sigma}_0^{-1} + \bm{\Sigma}_1^{-1})(\bm{\htau}_{\rm ST}-\bm{\tau})\le M\right)\\
  =& \P_{\bm{\Delta}}\left((\bm{\htau}_{\rm ST}-\bm{\tau})^\top (\bm{\Sigma}_0^{-1} + \bm{\Sigma}_1^{-1})(\bm{\htau}_{\rm ST}-\bm{\tau})\le M, \lVert (\bm{\Sigma}_0+ \bm{\Sigma}_1)^{-1/2}(\bm{\htau}_1 - \bm{\htau}_0) \rVert_2^2 \le q\right) \\
  &+ \P_{\bm{\Delta}}\left((\bm{\htau}_{\rm ST}-\bm{\tau})^\top (\bm{\Sigma}_0^{-1} + \bm{\Sigma}_1^{-1})(\bm{\htau}_{\rm ST}-\bm{\tau})\le M, \lVert (\bm{\Sigma}_0+ \bm{\Sigma}_1)^{-1/2}(\bm{\htau}_1 - \bm{\htau}_0) \rVert_2^2 > q\right)\\
  =& \Psi\left(M;\left\lVert (\bm{\Sigma}_0^{-1} + \bm{\Sigma}_1^{-1})^{-1/2} \bm{\Sigma}_1^{-1} \bm{\Delta}\right\rVert_2^2\right) \Psi\left(q; \left\lVert (\bm{\Sigma}_0 + \bm{\Sigma}_1)^{-1/2} \bm{\Delta}\right\rVert_2^2\right)\\
  &+ \int_{\lVert \bm{u} \rVert_2^2 > q} \Psi\left(M;\left\lVert (\bm{\Sigma}_0^{-1} + \bm{\Sigma}_1^{-1})^{-1/2} \bm{\Sigma}_1^{-1} [\bm{\Delta} - (1-h_q^*(\lVert \bm{u} \rVert_2^2))(\bm{\Sigma}_0+ \bm{\Sigma}_1)^{1/2} \bm{u}]\right\rVert_2^2\right) \phi\left(\bm{u};(\bm{\Sigma}_0 + \bm{\Sigma}_1)^{-1/2} \bm{\Delta}\right) \d \bm{u},
\end{align*}
which yields exactly the expression in \eqref{eq:P_multi_ST} by taking $\bm{\Sigma}^{-1/2}\bm{\Delta} = \bm{t}$.

Let $\bm{t}=\bm{\Sigma}^{-1/2}\bm{\Delta}$ and denote by
\[
p(\bm{t})
:=\P_{\bm{\Sigma}^{-1/2}\bm{\Delta}=\bm{t}}\!\left(
(\bm{\htau}_{\rm ST}-\bm{\tau})^\top(\bm{\Sigma}_0^{-1}+\bm{\Sigma}_1^{-1})
(\bm{\htau}_{\rm ST}-\bm{\tau})\le M
\right).
\]
By construction of the soft-thresholding rule,
$p(\bm{t})$ is invariant under coordinate-wise sign flips, i.e.,
$p(\bm{t})=p(\bm{t}\odot\bm{s})$ for all $\bm{s}\in\{\pm1\}^d$.
Moreover, for each $j\in\{1,\ldots,d\}$ and any fixed values of the remaining coordinates,
the map $u\mapsto p(t_1,\ldots,t_{j-1},u,t_{j+1},\ldots,t_d)$ is nonincreasing in $\lvert u \rvert$.
Therefore,
\[
\inf_{\bm{t}:|t_j|\le b_j} p(\bm{t})
=\inf_{\bm{s}\in\{\pm1\}^d} p(\bm{b}\odot\bm{s}),
\]
i.e., the worst-case coverage probability over the hyperrectangle is attained at a vertex.

Finally, choosing $\hat M_{\rm ST}$ to satisfy
\[
  \inf_{\bm{\Delta}: \bm{\Sigma}^{-1/2} \bm{\Delta} = \bm{b} \odot \bm{s}, \bm{s} \in \{-1,1\}^d} \P_{\bm{\Delta}}((\bm{\htau}_{\rm ST} - \bm{\tau})^\top (\bm{\Sigma}_0^{-1} + \bm{\Sigma}_1^{-1}) (\bm{\htau}_{\rm ST} - \bm{\tau}) \le M) = 1-\zeta
\]
completes the proof.

\end{proof}

\end{document}
